The standard model of cosmology, $\Lambda$CDM, provides a precise description of the universe and many of its defining features. As observational cosmology has entered a golden age of high-precision, some parameters have shown discrepancies that warrant extensive study and assessment, including the local rate of accelerated expansion, the Hubble constant ($H_0$). $H_0$ measurements derived from the early Universe -- most precisely, the \textit{Planck} cosmic microwave background analysis, infer $H_0 = 67.4\pm0.5\,\mathrm{km\,s^{-1}\,Mpc^{-1}}$ using the temperature power spectrum, E-mode polarization, and cross correlations \citep{Planck2020} - disagreeing with distance-ladder measurements from the local universe, most notably the SH0ES result of $H_0 = 73.0\pm1.0\,\mathrm{km\,s^{-1}\,Mpc^{-1}}$ \citep{Riess2022}. 
Despite extensive scrutiny of each approach, no systematic effect or new physics has been identified to resolve the discrepancy.
This $>5\sigma$ discrepancy is known as the Hubble Tension. 

\GW\ standard sirens offer an independent measurement of distances, and therefore, $H_0$. The \GW\ waveform from the coalescence of two compact objects directly encodes the luminosity distance to the source \citep{Schutz1986}. However, a \GW-based measurement alone cannot provide the source redshift, and several strategies have been developed to measure cosmological parameters using the \GW\ waveform information. Two strategies have been developed to combine \GW\ and electromagnetic data to enable a standard siren measurement, and are the focus of my research.

The first, the \textit{bright siren method}, requires the detection of an \EM\ counterpart to a \GW\ merger. The identification of a \GW\ merger candidate and the spectroscopic redshift of its host galaxy provides the highest $H_0$ precision per \GW\ event. The binary neutron star merger GW170817 remains the only \GW\ event with a confirmed \EM\ counterpart to date, and provided the first standard-siren measurement of $H_0$ \citep{2017ApJ...848L..12A,Abbott2017Nature,SoaresSantos2017}. The second method is the \textit{dark siren method}. Using dark sirens, it is possible to statistically marginalize over galaxies consistent with a \GW\ event's three-dimensional localization volume, weighting each galaxy by its probability of hosting the merger \citep{Gray2023}. This approach requires no \EM\ counterpart, and can be applied to all \GW\ detections coincident with galaxy catalog data.

I have designed a research program that will leverage standard sirens as an emerging probe of cosmology, maximizing the data from \LSST\ for bright and dark sirens, developing other astronomical instruments to leverage \LSST\ data, and ultimately using gravitational wave and electromagnetic data to address the Hubble Tension. 